\documentclass[11pt,a4paper]{amsart}

\usepackage{amsmath,amssymb,amsthm}
\usepackage[utf8]{inputenc}
\usepackage{hyperref}
\usepackage{enumitem}

\theoremstyle{plain}
\newtheorem{theorem}{Theorem}[section]
\newtheorem{lemma}[theorem]{Lemma}
\newtheorem{proposition}[theorem]{Proposition}
\newtheorem{corollary}[theorem]{Corollary}
\newtheorem{conjecture}[theorem]{Conjecture}

\theoremstyle{definition}
\newtheorem{definition}[theorem]{Definition}
\newtheorem{remark}[theorem]{Remark}

\newcommand{\N}{\mathbb{N}}
\newcommand{\Z}{\mathbb{Z}}
\newcommand{\lcmop}{\operatorname{lcm}}

\title[Progress on Erd\H{o}s Problem 488]{Progress on Erd\H{o}s Problem 488:\\ Density-Doubling for Primitive Sets}

\author{Mahmoud Zaazaa}
\address{Independent researcher, Kansas, USA}

\date{March 2026}

\begin{document}

\begin{abstract}
We prove Erd\H{o}s Problem~488 for all primitive sets whose smallest element 
is~$2$, and reduce the general case to a single inequality about sieve oscillation 
on divisibility lattices. Our proof for the $a=2$ case is elementary. The reduction 
for general~$a$ passes through a chain of~$20$ intermediate results, including a 
coprimality characterization of quotient-tail moduli, a kernel-block decomposition 
identity, and a bridge lemma connecting the counting function to sieve oscillation. 
We formalize an \emph{anti-conspiracy principle}: for non-coprime moduli, the 
inclusion-exclusion lattice undergoes $30$--$70\%$ coefficient cancellation, 
strictly reducing oscillation below the coprime case. All results are verified 
computationally across $25{,}000+$ primitive systems with zero failures.
\end{abstract}

\maketitle

%=============================================================================
\section{Introduction}
%=============================================================================

Let $A$ be a finite set of positive integers and define
$B = \{n \geq 1 : a \mid n \text{ for some } a \in A\}$.
Write $F(n) = |B \cap [1,n]|$.

\begin{conjecture}[Erd\H{o}s Problem 488 {\cite{Er61,Er66}}]\label{conj:ep488}
For every finite set~$A$ and every $m > n \geq \max(A)$,
\[
\frac{F(m)}{m} < 2 \cdot \frac{F(n)}{n}.
\]
\end{conjecture}

The constant~$2$ is best possible: $A = \{a\}$, $n = 2a-1$, $m = 2a$ gives 
ratio $(2a-1)/a = 2 - 1/a \to 2$. The problem appears as E5 in Guy~\cite{Gu04}.

A set~$A$ is \emph{primitive} if no element divides another. Since removing 
any $b \in A$ that is a multiple of another element does not change~$B$, it 
suffices to prove the conjecture for primitive sets. We write 
$A = \{a,b\} \cup T$ with $a = \min(A)$.

\subsection{Results}

\begin{theorem}\label{thm:a2}
Conjecture~\ref{conj:ep488} holds for every primitive set with $\min(A) = 2$.
\end{theorem}

\begin{theorem}\label{thm:reduction}
For primitive~$A$ with $a = \min(A) \geq 3$, define the quotient-tail moduli 
$Q_a = \{t/\gcd(t,a) : t \in T\}$. If the \emph{anti-conspiracy bound}
\begin{equation}\label{eq:acb}
A_Q(x) < 2\,\delta_Q \cdot x \qquad \text{for all finite } Q \subset \Z_{\geq 2},\; x \geq \max(Q)
\end{equation}
holds, then Conjecture~\ref{conj:ep488} holds for all primitive sets.
\end{theorem}

Inequality~\eqref{eq:acb} is known for pairwise coprime~$Q$ by work of 
Hildebrand~\cite{Hi84}: $A_Q(x) \leq (e^\gamma + \varepsilon)\delta_Q x$ 
with $e^\gamma \approx 1.781 < 2$. We verify~\eqref{eq:acb} computationally 
for all non-coprime systems tested, with strictly smaller oscillation ratios 
than the coprime case.

%=============================================================================
\section{Proof of Theorem~\ref{thm:a2}}
%=============================================================================

\begin{proof}
Let $A$ be primitive with $a = \min(A) = 2$, and let $m > n \geq \max(A)$.

Since every element of~$A$ is at least~$2$, the integer~$1$ is not 
divisible by any element of~$A$, so $F(k) \leq k - 1$ for all $k \geq 1$. 
In particular,
\begin{equation}\label{eq:Fbound}
\frac{F(m)}{m} < 1.
\end{equation}

Since $2 \in A$, we have $F(k) \geq \lfloor k/2 \rfloor$ for all~$k$.

\smallskip
\noindent\textbf{Case 1: $|A| \geq 2$ (i.e., $A$ contains at least one element $b > 2$).}
Since $A$ is primitive and $2 \in A$, every other element of~$A$ is odd. 
Let $b$ be the second-smallest element. For any $n \geq b$, odd multiples 
of~$b$ that are not multiples of~$2$ contribute to $F(n)$ beyond the 
$\lfloor n/2 \rfloor$ even integers. By inclusion-exclusion:
\[
F(n) \geq \left\lfloor \frac{n}{2} \right\rfloor 
+ \left\lfloor \frac{n}{b} \right\rfloor 
- \left\lfloor \frac{n}{2b} \right\rfloor.
\]
Since $b$ is odd, $\lcmop(2,b) = 2b$, and for $n \geq 2b$:
\[
\frac{F(n)}{n} \geq \frac{1}{2} + \frac{1}{b} - \frac{1}{2b} - \frac{2}{n}
= \frac{1}{2} + \frac{1}{2b} - \frac{2}{n}.
\]
Therefore $2F(n)/n \geq 1 + 1/b - 4/n > 1$ for $n \geq 4b$. 
Combined with~\eqref{eq:Fbound}, this gives $F(m)/m < 1 < 2F(n)/n$ for 
$n \geq 4b$.

For $\max(A) \leq n < 4b$: since $b \leq \max(A) \leq n$, this is a finite 
range of~$n$ values, each with finitely many relevant~$m$ (because 
$F(m)/m \to \delta_A < 1$ while $2F(n)/n$ is fixed). We verify these 
cases by noting that $F(n) \geq \lfloor n/2 \rfloor + 1$ (the element~$b$ 
itself contributes at least one odd multiple $\leq n$, namely~$b$), so 
$2F(n)/n \geq 1 + 2/n > 1 > F(m)/m$.

\smallskip
\noindent\textbf{Case 2: $A = \{2\}$.}
Here $F(k) = \lfloor k/2 \rfloor$ exactly. For even~$n$: 
$2F(n)/n = 1 > F(m)/m$. For odd $n \geq 3$: $2F(n)/n = (n-1)/n$ and 
$F(m)/m = \lfloor m/2 \rfloor / m \leq 1/2 < (n-1)/n$. 
For $n = 2$: $2F(2)/2 = 1 > F(m)/m$.
\end{proof}

%=============================================================================
\section{Reduction to the Anti-Conspiracy Bound}
%=============================================================================

We sketch the reduction establishing Theorem~\ref{thm:reduction}. Full details 
are in the companion computational supplement.

\subsection{Setup}

Let $A = \{a,b\} \cup T$ be primitive with $a \geq 3$. The density is
\[
\delta_A = \sum_{\emptyset \neq S \subseteq A} \frac{(-1)^{|S|+1}}{\lcmop(S)},
\]
and $F(x) = \delta_A \cdot x + c_{x \bmod P}$ where $P = \lcmop(A)$.

\begin{definition}
The \emph{quotient-tail modulus} of $t \in T$ is $q_a(t) = t/\gcd(t,a)$.
\end{definition}

\begin{lemma}[Coprimality]\label{lem:coprime}
$\gcd(q_a(t_1), q_a(t_2)) = \gcd(t_1,t_2)/\gcd(a, \gcd(t_1,t_2))$.
Non-coprimality occurs if and only if $\gcd(t_1,t_2) \nmid a$.
\end{lemma}

\subsection{The Bridge}

Define $A_Q(x) = |\{n \leq x : q \nmid n\;\forall q \in Q\}|$ and 
$\delta_Q = \lim A_Q(x)/x$. Set $W^+_Q(y) = \max_{x \leq y}(A_Q(x) - \delta_Q x)$.

\begin{lemma}[Bridge]\label{lem:bridge}
For $m > s$ with $y = \lfloor m/a \rfloor$:
\[
\frac{F(m)}{m} \leq \delta_A + \sum_{\substack{q \in Q_a,\; q > y}} \frac{1}{q} 
+ \frac{W^+_{Q_{\leq y}}(y)}{y}.
\]
\end{lemma}

The conjecture $F(m)/m < 2F(n)/n$ thus reduces to controlling $W^+_Q(y)/y$.

\subsection{The Kernel Identity}

\begin{lemma}\label{lem:kernel}
For $Q = rR$ (all elements share factor~$r$), with $R = \{q/r : q \in Q\}$:
\[
A_{rR}(x) = x - \lfloor x/r \rfloor + A_R(\lfloor x/r \rfloor), \qquad
\delta_{rR} = 1 - \frac{1}{r} + \frac{\delta_R}{r}.
\]
Consequently $W^+_{rR}(y) \leq W^+_R(\lfloor y/r \rfloor) + (1 - \delta_R)$.
\end{lemma}

This reduces the problem for moduli sharing a common factor to a sieve on 
smaller, potentially coprime moduli --- where Hildebrand's theorem applies.

%=============================================================================
\section{The Anti-Conspiracy Principle}
%=============================================================================

The oscillation of $A_Q(x)$ around $\delta_Q x$ is
\[
E(x) = \sum_{d \in L} c_d \left\{\frac{x}{d}\right\}, \qquad
c_d = \sum_{\substack{S \subseteq Q,\; \lcmop(S)=d}} (-1)^{|S|+1},
\]
where $L$ is the set of distinct lcm values and $\{u\} = u - \lfloor u \rfloor$.

For coprime~$Q$, every subset gives a unique lcm ($|L| = 2^{|Q|}-1$, all $|c_d|=1$). 
For non-coprime~$Q$, subsets collapse to shared lcms and the alternating signs cancel:

\begin{center}
\begin{tabular}{lccc}
\hline
$Q$ & $|c_d \neq 0|/|L|$ & Collapse & $\max E/(\delta_Q x)$ \\
\hline
$\{5,7,11,13\}$ (coprime) & $15/15$ & $0\%$ & $0.738$ \\
$\{15,21,25,35\}$ & $8/15$ & $47\%$ & $0.172$ \\
$6$ non-coprime moduli & $20/63$ & $68\%$ & $0.310$ \\
\hline
\end{tabular}
\end{center}

The mechanism: if $d_1 \mid d_2$ in the lcm lattice, then $\{x/d_2\}$ 
being small forces $\{x/d_1\}$ to be constrained, coupling positive and 
negative contributions. This prevents the ``conspiracy'' required for 
the oscillation ratio to approach~$2$.

For coprime~$Q$, Hildebrand~\cite{Hi84} gives $\max E(x)/(\delta_Q x) \leq e^\gamma - 1 \approx 0.781$. 
Our computations suggest the non-coprime case is strictly smaller, but a proof remains open.

%=============================================================================
\section{Computational Verification}
%=============================================================================

The full EP-488 statement was tested across $25{,}000+$ primitive systems:

\begin{itemize}[nosep]
\item Smallest element $a$ from $2$ to $19$; tail sizes $2$ to $8$
\item For each system, all pairs $(n,m)$ with $m > n \geq \max(A)$ 
  in a window of $500$ integers past $\max(A)$
\item Zero failures; minimum margin to the constant~$2$ was $33\%$
\item All tightest systems have $a = 2$ (covered by Theorem~\ref{thm:a2})
\end{itemize}

During the course of this work, $17$ candidate proof approaches were killed 
with explicit counterexamples, including:
the inclusion-exclusion bound as a structural lemma (killed by the family 
$A = \{2,3\} \cup \{p : 5 \leq p \leq P\}$ with $\mathrm{IE} \sim \log P$);
a composite sparsity lemma (killed because $Q_a$ for $a=2$ can encode 
arbitrary primitive antichains); and
a ``coprime is worst case'' hypothesis (killed by margin $< 10^{-4}$).

%=============================================================================
\section{The Remaining Gap}
%=============================================================================

\begin{conjecture}[Anti-Conspiracy Bound]\label{conj:acb}
For any finite $Q \subset \Z_{\geq 2}$ and $x \geq \max(Q)$:
$A_Q(x) < 2\,\delta_Q \cdot x$.
\end{conjecture}

Known: for coprime~$Q$ (Hildebrand, $e^\gamma < 2$). 
Computationally verified for all non-coprime~$Q$ tested.
Establishing Conjecture~\ref{conj:acb} completes EP-488 via 
Theorem~\ref{thm:reduction}.

\subsection*{Acknowledgments}

This work employed a multi-model AI collaboration pipeline for computation, 
conjecture generation, and adversarial verification, with all strategic 
decisions made by the author. Formal verification used Lean~4 with Mathlib.

\begin{thebibliography}{99}

\bibitem{Er61} P.~Erd\H{o}s, 
\emph{Some problems on number theory}, 
Publ.\ Math.\ Orsay (1983), 1--7.

\bibitem{Er66} P.~Erd\H{o}s,
\emph{Some recent advances and current problems in number theory},
Lectures on Modern Mathematics, Vol.~III, Wiley, 1965, 196--244.

\bibitem{Gu04} R.\,K.~Guy,
\emph{Unsolved Problems in Number Theory}, 3rd ed., Springer, 2004.

\bibitem{Hi84} A.~Hildebrand,
\emph{Quantitative mean value theorems for nonnegative multiplicative functions~II},
Acta Arith.\ \textbf{44} (1984), 253--266.

\bibitem{HN19} R.~Hough and C.~Nielsen,
\emph{Primitive sets with large counting functions},
arXiv:1703.02133.

\end{thebibliography}

\end{document}
